\section{Conclusion}
\label{sec:conclusion}

Two coding agents, given an identical checkout of GraphViz and a prompt naming
no target, both located the dominant cost in the \texttt{dot} layout pipeline
and produced output-preserving speedups: $1.47\times$ and $1.16\times$
end-to-end at 5{,}000 nodes, against an $8.4\%$ measured noise floor. Finding
the bottleneck was evidently not the difficult part of the task.

What separated them was scope. At 10{,}000 nodes the dominant phase shifts from
crossing minimization to coordinate assignment, and the agent that had
restructured network simplex --- code shared by both phases --- improves to
$1.57\times$ while the agent that optimized a single function stays flat at
$1.15\times$. An optimization's value depends on whether the phase it targets
grows or shrinks with the workload, and a single problem size cannot reveal
that.

Neither agent used parallelism. On this workload the ceiling arithmetic suggests
that was the right call: at 10{,}000 nodes, idealized eight-core parallelization
of crossing minimization caps end-to-end speedup at $1.49\times$, and one agent
exceeded that serially at $1.57\times$ --- without introducing concurrency or the
possibility of a race.

The results are modest, and we have tried to report them at the size they
actually are. The protocol may be the more durable contribution: measure the
noise floor with byte-identical binaries before measuring anything else, hold
the evaluation harness out of the agents' reach, forbid the shortcuts that
manufacture speedups, and attribute results by phase at more than one problem
size. Each of these caught something in this study that a simpler design would
have reported as a finding.

Immediate future work is to repeat both agent runs to establish whether the
ordering is stable, extend to a third graph size and to non-layered topologies,
and complete the correctness checks named in Section~\ref{sec:limitations}.
